\documentclass[../../../include/open-logic-section]{subfiles}

\begin{document}

\olfileid[pt]{sth}{choice}{tarskiscott}
\olsection{O Truque de Tarski--Scott}

Na \olref[cardinals][cardsasords]{defcardinalasordinal},
definimos cardinais como ordinais. Para isso, assumimos o Axioma da
Boa Ordenação. Fizemo-lo apenas porque é o
``padrão da indústria''.

Antes de discutir questões filosóficas em torno da
Boa Ordenação, importa deixar claro que \emph{podemos}
afastar-nos desse padrão e desenvolver uma teoria de cardinais
\emph{sem} assumir a Boa Ordenação. Continuamos a poder usar as
definições de $\cardeq{A}{B}$, $\cardle{A}{B}$ e $\cardless{A}{B}$,
tal como apareceram no \olref[sfr][siz][]{chap}. Só precisamos de uma nova
noção de \emph{cardinal}.

Uma ideia ingênua seria tentar definir a cardinalidade de $A$ assim:
\[
	\Setabs{x}{\cardeq{A}{x}}.
\]
Talvez queira comparar isto com a definição fregeana de $\# x Fx$,
esboçada no final de \olref[cardinals][hp]{sec}. E, pelas razões
que ali adumbramos, esta definição falha. Qualquer conjunto singular
é equipotente com $\{\emptyset\}$. Mas novos singulares são
formados em cada estágio sucessor da hierarquia (basta considerar o
singular do estágio anterior). Logo $\Setabs{x}{\cardeq{A}{x}}$ não
existe, pois não pode ter posto.

Para contornar o problema, usamos um truque devido a Tarski e Scott:\footnote{Lembrete: todas as fórmulas podem ter parâmetros (salvo indicação em contrário).}

\begin{defn}[Tarski--Scott]
Para qualquer fórmula $\phi(x)$, seja
$[ x : \phi(x)] $ o conjunto de todos os $x$, de posto mínimo possível, tais que
$\phi(x)$ (ou $\emptyset$, se não houver $\phi$s).
\end{defn}

Devemos verificar que esta definição é legítima. Trabalhando em $\ZF$,
\olref[spine][foundation]{zfentailsregularity} garante que
$\setrank{x}$ existe para todo o $x$. Agora, se houver entidades
que satisfazem $\phi$, podemos deixar $\alpha$ ser o posto mínimo tal
que  $(\exists x\subseteq V_\alpha)\phi(x)$, O.s.\@, $(\forall \beta
\in \alpha)(\forall x \subseteq V_\beta)\lnot \phi(x)$. Podemos então
definir $[x : \phi(x)]$ por Separação como $\Setabs{x \in
V_{\alpha+1}}{\phi(x)}$. 

Tendo justificado o truque de Tarski--Scott, podemos usá-lo para definir
uma noção de cardinalidade:

\begin{defn}
A \textsc{ts}-cardinalidade de $A$ é $\text{tsc}(A) = [x :
\cardeq{A}{x}]$.
\end{defn}

A definição de \textsc{ts}-cardinal não usa a Boa Ordenação.
Mas, mesmo sem esse Axioma, podemos mostrar que os
\emph{\textsc{ts}-cardinais} se comportam de modo semelhante aos \emph{cardinais} tal como
definidos na \olref[cardinals][cardsasords]{defcardinalasordinal}.
Por exemplo, se reformularmos
\olref[cardinals][cardsasords]{lem:CardinalsBehaveRight} e
\olref[card-arithmetic][opps]{lem:SizePowerset2Exp} em termos de
\textsc{ts}-cardinais, as provas funcionam em $\ZF$,
sem assumir a Boa Ordenação. 

Enquanto estamos no assunto, vale a pena notar que também podemos
desenvolver uma teoria de ordinais usando o truque de Tarski--Scott. Onde
$\tuple{A, <}$ é uma boa ordenação, seja $\text{tso}(A, <) = [\tuple{X,
R} : \ordeq{\tuple{A, <}}{\tuple{X, R}}]$. Para mais sobre este tratamento
de cardinais e ordinais, ver \citet[chs.~9--12]{Potter2004}.

\end{document}
